% Regression: a dot sits at the middle of the rectangle its atom claims, and
% the label band is shorter than one pitch.  Across an EVEN span the dot
% stands halfway between two lanes, so the two stations facing them land
% inside those lanes rather than clear of them; across an odd span the dot
% stands on a lane of its own and every station falls in the daylight
% between lanes.  A wire running the length of a lane an even span stations
% in therefore takes that face, and not only the face it leaves by.
\documentclass{standalone}
\usepackage{tenkz}
\makeatletter
\ExplSyntaxOn
\file_input:n { tenkz-kernel.code.tex }
\int_new:N \l__tenkz_test_side_count_int
\tl_new:N \l__tenkz_test_side_label_tl
\cs_new_eq:NN
  \__tenkz_test_side_auto:nN \__tenkz_kernel_r_label_auto_side:nN
\cs_set_protected:Npn \__tenkz_kernel_r_label_auto_side:nN #1#2
  {
    \__tenkz_test_side_auto:nN {#1} #2
    \int_gincr:N \l__tenkz_test_side_count_int
    \__tenkz_model_get:nnN {#1} {label} \l__tenkz_test_side_label_tl
    \str_case:VnF \l__tenkz_test_side_label_tl
      {
        {L}
          {
            \str_if_eq:VnF #2 {w}
              { \errmessage{a~two-row~span~stationed~north~in~its~own~row} }
          }
        {C}
          {
            \str_if_eq:VnF #2 {w}
              { \errmessage{a~four-column~span~stationed~east~in~its~own~column} }
          }
        {O}
          {
            \str_if_eq:VnF #2 {n}
              { \errmessage{an~odd~span~gave~up~a~station~standing~in~daylight} }
          }
      }
      { \errmessage{unexpected~label~in~span-lane~fixture} }
  }
\ExplSyntaxOff
\makeatother
\begin{document}
\tenkzkernel
% Two rows: the north station lands in row 1, where the wire east runs, so
% only west is left.
\begin{tenkz}[lattice={3x2}, bonds=none]
  \tn[at=(1,1), wires=2, skin=dot]{L}
  \tn[at=(1,2)]{}
  \tn[at=(3,1)]{}
  \tnwire{(1,1)}{(1,2)}
  \tnwire{(2,1)}{(3,1)}
\end{tenkz}
% Four columns: the east station lands in column 3, an interior column of
% the span, where the wire south runs.  The wires on column 1 take north and
% south without touching a lane.
\begin{tenkz}[lattice={3x5}, bonds=none]
  \tn[at=(2,1), wide=4, skin=dot]{C}
  \tn[at=(1,1)]{}
  \tn[at=(3,1)]{}
  \tn[at=(3,3)]{}
  \tnwire{(2,1)}{(1,1)}
  \tnwire{(2,1)}{(3,1)}
  \tnwire{(2,3)}{(3,3)}
\end{tenkz}
% Three rows: the dot stands on row 2 and its north station falls between
% rows, clear of the wire running east along row 1, so north stays free.
\begin{tenkz}[lattice={4x2}, bonds=none]
  \tn[at=(1,1), wires=3, skin=dot]{O}
  \tn[at=(1,2)]{}
  \tn[at=(4,1)]{}
  \tnwire{(1,1)}{(1,2)}
  \tnwire{(3,1)}{(4,1)}
\end{tenkz}
\ExplSyntaxOn
\int_compare:nNnF { \l__tenkz_test_side_count_int } = {3}
  { \errmessage{span-lane~fixture~did~not~station~three~labels} }
\ExplSyntaxOff
\end{document}
